\documentclass[10pt,a4paper]{amsart}
\usepackage[margin=19mm]{geometry}
\usepackage{amsmath,amssymb,mathtools}
\usepackage[scr=rsfs,cal=euler]{mathalfa}
\usepackage{xfrac}
\usepackage[unicode,hidelinks,bookmarks=false]{hyperref}
\newcommand{\ie}{i.e.\ }
\newcommand{\cf}{cf.\ }
\newcommand{\resp}{respectively\ }
\newcommand{\Subsection}{\subsection}
\providecommand{\nobreakdash}{\nobreak\hbox{-}}
\setlength{\emergencystretch}{2em}
\multlinegap0pt
\usepackage{xcolor}
\usepackage{enumerate}
\usepackage{float}
\usepackage{amssymb}
\usepackage{breqn}
\usepackage{mathtools}
\usepackage{tikz}
\newtheorem{proposition}{Proposition}
\newtheorem{theorem}{Theorem}
\newcommand{\csm}{\textrm{csm}}
\newcommand{\perm}{\textrm{perm}}
\title{Three formulas for CSM classes of open quiver loci}
\author{Moriah Elkin}
\date{Source: arXiv:2511.17723v4 (CC BY 4.0)}
\begin{document}
\maketitle

\noindent\emph{Source and licence.} Three formulas for CSM classes of open quiver loci. Version arXiv:2511.17723v4 (CC BY 4.0), distributed by arXiv under the Creative Commons Attribution 4.0 International licence, http://creativecommons.org/licenses/by/4.0/, as recorded in the arXiv licence metadata for this exact version and shown on the abstract page https://arxiv.org/abs/2511.17723v4. Full licence notice: \texttt{licenses/arxiv-2511.17723-CC-BY-4.0.txt}.\par\smallskip

\noindent\textit{Editorial scope.} Counted unit: the article's theorem thm:cgpdopen (source0.tex lines 1589--1592). The article's complete original proof of it is delivered with it (source0.tex lines 1670--1675, one \texttt{proof} environment of 6 lines, which the article places away from the statement). No mathematics is written, repaired or re-derived: the statement and the proof are the article's own lines, in the article's own notation, and no retained line is substituted or re-flowed.  Retained uncounted as the source-local context the proof needs: lines 1368--1374.  Nothing was omitted from any retained span, so the package carries no omission record.  No retained line contains a citation, so the wrapper carries no bibliography and no bibliography entry.  No retained line contains a figure environment, an image-inclusion command or drawing code, so no image is delivered and the package declares no asset dependency.  Editorial formatting only, in addition: an AMS \texttt{amsart} preamble that restates the article's own declarations and macros used by the retained lines, namely the environments \texttt{proposition}, \texttt{theorem} on separate counters, together with the standard packages those lines need.  The retained environments print in this document as Proposition 1, Theorem 1 in order of appearance, whereas the article prints the same items with the numbering of its own full document.  The complete native source of the article as distributed in the arXiv e-print for this exact version is bundled unchanged as \texttt{sources/189606/source0.tex}.
\begin{proposition}[Pipe dream formula for open quiver loci]\label{prop:pdopen}
For any $v \in S_d$, let $\mathcal{P}(v)$ be the set of (not necessarily reduced!) pipe dreams for $v$, and let $c(D)$ be the set of cross tiles in a given pipe dream $D$, so $|c(D)|$ is the number of cross tiles in $D$.
Given a rank array $\mathbf{r}$, we have the following expression for the CSM class of the corresponding open quiver locus:
\[\csm(\Omega_{\mathbf{r}}^\circ \subset P_- \backslash GL_d) =
\sum_{v \in \perm(\mathbf{r})} 
\sum_{D \in \mathcal{P}(v)} \hbar^{\ell(w_0\mathbf{w}_0)-|c(D)|}(\mathbf{x}-\mathring{\mathbf{x}})^{D\setminus D_{Hom}}.\]
\end{proposition}

\begin{theorem}[CGPD formula for open quiver loci]\label{thm:cgpdopen}
Given a rank array $\mathbf{r}$, we have the following expression for the CSM class of the corresponding open quiver locus:
\[\csm(\Omega_{\mathbf{r}}^\circ \subset Hom) = \sum_{\delta \in \mathcal{CGPD}(\mathbf{r})} (\mathbf{x}-\mathring{\mathbf{x}})^\delta.\]
\end{theorem}

\begin{proof}[Proof of Theorem \ref{thm:cgpdopen}]
This is a purely combinatorial translation of Proposition \ref{prop:pdopen}. We first note that all tiles except those on the block superantidiagonal are completely determined and may be deleted, replacing the bumps on the block diagonal with pipes continuing due southwest. There are $\ell(w_0\mathbf{w}_0)-|c(D)|$ remaining bump tiles, so we give these tiles a weight of $\hbar$. Next, note that the conditions on where pipes enter and exit the rectangles exactly encode whether the pipe dream's permutation $v$ is an element of $\perm(\mathbf{r})$. The lace array $\mathbf{s}$ is reproduced on and below the block antidiagonal of $z(\mathbf{r})$, and if the matrix of $z(\mathbf{r})$ has a 1 in block row $i$ and block column $j$ for $i \leq j$, then a pipe must enter from the right in rectangle $i$ and leave from the left in rectangle $j$.
Not requiring specific entrance or exit locations \emph{within} a particular rectangle corresponds to varying the locations of the rank jumps of $z(\mathbf{r})$ \emph{within} blocks to form $v \in \perm(\mathbf{r})$.

In fact, since we only care about what the pipes representing laces do, we can erase the other pipes; if a tile is completely blank, the erased pipes may cross or not, so its weight is the sum of a bump tile and a cross tile. If two pipes of the same color meet, whether they cross or not doesn't affect the condition from \textbf{r}, so the weights of both pipe dreams should always be included. We can therefore arbitrarily disallow same-color pipes from crossing in the diagrams while assigning the same-color bump tile the sum of a bump and a cross. This brings us to the definition of CGPDs.
\end{proof}

\end{document}
